\documentclass[11pt]{article}
\usepackage[margin=1in]{geometry}
\usepackage{array}
\usepackage{booktabs}
\usepackage{enumitem}
\usepackage{hyperref}
\usepackage{xcolor}
\usepackage{longtable}
\usepackage{xurl}
\usepackage{fvextra}
\usepackage[T1]{fontenc}
\usepackage{lmodern}

\hypersetup{colorlinks=true, linkcolor=black, urlcolor=blue}
\setlist[itemize]{leftmargin=1.2em, itemsep=0.2em, topsep=0.2em}
\setlist[enumerate]{leftmargin=1.4em, itemsep=0.2em, topsep=0.2em}
\emergencystretch=3em
\newcommand{\code}[1]{\texttt{#1}}
\newcommand{\status}[1]{\textbf{\texttt{#1}}}
\DefineVerbatimEnvironment{verbatim}{Verbatim}{breaklines=true,breakanywhere=true}

\title{ARENA 3.0 Extended Review Packet: DiffusionGemma NVFP4 Proof and BF16 Deferral}
\author{Local verification packet}
\date{2026-07-01}

\begin{document}
\maketitle

\section*{Verdict}
The section 5.5 diffusion-language-model exercises currently have a valid CUDA
toy-model learning ladder, deterministic controls, and a strict claim boundary.
They now also have a real released-checkpoint DiffusionGemma generation proof
for the 24GB route: the pinned NVIDIA NVFP4 checkpoint generated a non-empty
answer in an isolated vLLM 0.24.0 runtime on the local RTX 5090 Laptop GPU. The
main uv environment remains the requested current stack, torch 2.12.1+cu132 with
CUDA 13.2, and vLLM is not installed into it because public vLLM 0.24.0 resolves
to torch 2.11.0+cu130. The full Google BF16 DiffusionGemma path remains
\status{deferred} for direct local loading because the required weight shards
alone are 51,647,701,024 bytes, or 48.10 GiB, while the machine reports 23.46
GiB of GPU memory.

\section{Motivation and Research Question}
The motivation is to keep the extension faithful to original ARENA standards:
students should implement the core algorithm, test it against explicit
references and controls, and only then connect it to a real modern model. The
research question is: can this extension honestly claim an ARENA-style
DiffusionGemma exercise on a single 24GB laptop GPU?

The current answer is narrower:
\begin{itemize}
  \item Yes for the pedagogical discrete-diffusion exercises: the implementation
  has noising, denoising loss, remasking, sampler, activation-trajectory checks,
  CUDA training, and a shuffled-label negative control.
  \item Yes for released-checkpoint generation in the intended 24GB route:
  NVIDIA NVFP4 has a pinned isolated-vLLM proof artifact with prompt, output,
  runtime, and GPU metadata.
  \item Deferred for direct Google BF16 local inference on this machine: the model
  weights alone exceed the local VRAM target.
  \item Still not claimed: DiffusionGemma denoising-step activation capture,
  diffusion-time patching, quality evaluation, throughput benchmarking, or full
  interpretability parity with original ARENA-style mechanistic investigations.
\end{itemize}

\section{Background}
DiffusionGemma is a text diffusion language model. Unlike a standard
autoregressive transformer that commits to one next token at a time, a discrete
diffusion language model starts with masked or noisy tokens and iteratively
denoises a block. The course exercise teaches this by first implementing a small
block denoising model on generated integer-token data, then checking whether the
released DiffusionGemma artifacts can be used locally.

The local hardware target for this fork is a single 24GB CUDA GPU. The current
machine reports:
\begin{itemize}
  \item GPU: NVIDIA GeForce RTX 5090 Laptop GPU.
  \item GPU memory reported by torch: 23.46 GiB.
  \item uv-managed torch stack: torch 2.12.1+cu132, CUDA 13.2, torchvision
  0.27.1+cu132.
  \item ModelOpt: nvidia-modelopt 0.44.0 is installed.
  \item Isolated vLLM proof runtime: vLLM 0.24.0, torch 2.11.0+cu130, CUDA 13.0.
\end{itemize}

\section{Dataset and Exercise Structure}
The accepted part of section 5.5 uses a generated conditional copy-pair grammar.
Each example contains two protected prefix tokens and four suffix tokens. The
label rule is:
\[
  [a,b] \mapsto [a,a,b,b].
\]
For example, the protected prefix \([3,5]\) gives the full sequence
\([3,5,3,3,5,5]\). The verification report records 80 training examples, 20
held-out examples, vocabulary size 11, sequence length 6, mask token id 10, and a
deterministic shuffled-label control.

\section{Methodology}
The toy exercise checks the following mathematical objects:
\begin{enumerate}
  \item A linear masking schedule \(p_t\) that controls how much of the suffix is
  masked at denoising step \(t\).
  \item A masked-token denoising objective:
  \[
    L = -\frac{1}{|M|}\sum_{i \in M}\log p_\theta(x_i \mid x_{\setminus M}, t),
  \]
  where \(M\) is the set of masked positions.
  \item A confidence-remasking sampler that repeatedly predicts suffix tokens and
  re-masks low-confidence positions.
  \item A commitment-time trace that records when each suffix position becomes
  stable during denoising.
\end{enumerate}
This tests the teaching mechanism, not the released checkpoint. The released
checkpoint path is checked separately by exact repository revisions, required
files, shard counts, runtime metadata, and an explicit boolean
\code{diffusiongemma\_generation\_ready}.

\subsection*{Pseudocode}
\begin{verbatim}
for batch in generated_copy_pair_examples:
    t = sample_diffusion_step()
    masked = apply_linear_suffix_mask(batch.tokens, t)
    logits = denoiser(masked.tokens, t)
    loss = cross_entropy(logits[masked.positions], batch.tokens[masked.positions])
    update_model(loss)

for heldout_batch in heldout_examples:
    sample = confidence_remasking_sampler(denoiser, heldout_batch.prefix)
    record_exact_match_suffix_accuracy(sample, heldout_batch.target)
    record_commitment_times_and_entropy(sample.trajectory)

run_same_training_path_with_shuffled_suffix_labels()
assert shuffled_label_accuracy_is_low()
assert released_diffusiongemma_generation_ready
assert nvfp4_generation_proof_runtime_is_isolated_from_main_uv
assert bf16_direct_local_loading_is_deferred_for_24gb
\end{verbatim}

\section{Metrics Primer and Results}
The primary toy metrics are held-out masked-token accuracy, sampler exact-match
rate, shuffled-label control accuracy, and measured peak VRAM. Held-out accuracy
checks whether the denoiser learned the generated copy-pair rule. Sampler
exact-match checks whether iterative denoising reconstructs the entire suffix.
The shuffled-label control should fail, showing the model is not passing through
an accidental shortcut. Matthews correlation coefficient (MCC) is not applicable
because this is not a binary classification task.

\begin{center}
\begin{tabular}{@{}ll@{}}
\toprule
Metric & Current observed value \\
\midrule
Section 5.5 verification accepted & true \\
Toy CUDA held-out masked accuracy & 1.0 \\
Toy CUDA sampler exact match & 1.0 \\
Shuffled-label control accuracy & 0.0714 \\
Toy CUDA peak VRAM & 0.0737 GiB \\
DiffusionGemma released-checkpoint generation ready & true \\
NVFP4 isolated vLLM proof & true, vLLM 0.24.0 / torch 2.11.0+cu130 \\
Main uv vLLM generation support & false, main uv stays torch 2.12.1+cu132 \\
Google BF16 weight bytes required & 51,647,701,024 bytes, 48.10 GiB \\
NVIDIA NVFP4 weight bytes required & 18,823,855,888 bytes, 17.53 GiB \\
\bottomrule
\end{tabular}
\end{center}

\section{Qualitative Examples}
The toy qualitative example is a generated grammar row: prefix \([3,5]\) implies
target suffix \([3,3,5,5]\), so the full sequence is \([3,5,3,3,5,5]\). Passing
toy samples should reconstruct this suffix from masked states, while the
shuffled-label control should not learn a stable rule.

The released-checkpoint proof used this prompt:
\begin{quote}
\small In one concise sentence, explain why mechanistic interpretability
experiments need negative controls.
\end{quote}
The pinned NVFP4 checkpoint generated:
\begin{quote}
\small Negative controls are necessary to ensure that observed model behaviors
are causally linked to the specific mechanisms being tested rather than
artifacts, chance correlations, or noise.
\end{quote}
This is a one-prompt runtime proof, not an output-quality benchmark.

\section{Artifact Status}
\begin{longtable}{@{}p{0.20\linewidth}p{0.17\linewidth}p{0.55\linewidth}@{}}
\toprule
Item & Status & Evidence and implication \\
\midrule
\endhead
Google BF16 DiffusionGemma &
\status{deferred} &
\code{google/diffusiongemma-26B-A4B-it} at revision
\code{0f28bc42f588fbd8f71e08102b1c3960298a1358}. The report requires 11
safetensor shards totaling 48.10 GiB. Local shards are incomplete. This is too
large for direct local loading on a 23.46 GiB GPU, before activations and runtime
overhead. \\
\addlinespace
NVIDIA NVFP4 DiffusionGemma &
\status{proven in isolated runtime} &
\code{nvidia/diffusiongemma-26B-A4B-it-NVFP4} at revision
\code{2ea837236295d617ac27f8c17d61228081932c40}. The two weight shards total
17.53 GiB and are the intended 24GB path. The local shards are complete, the
model path matches the pinned revision, and
\path{artifacts/diffusiongemma_vllm_probe.json} records a real generation on the
23.46 GiB RTX 5090 Laptop GPU. \\
\addlinespace
Transformers NVFP4 loading &
\status{not the active proof path} &
The checkpoint advertises \code{quant\_method: modelopt}. The current
Transformers quantization registry rejects this with ``Unknown quantization type,
got modelopt''. Config, tokenizer, processor, and model-class support remain
useful preflight evidence, but the generation proof currently comes from vLLM,
not Transformers direct loading. \\
\addlinespace
vLLM dependency path &
\status{isolated} &
\code{uv pip install --dry-run vllm --prerelease allow} resolves vLLM 0.24.0,
but would replace torch 2.12.1+cu132 with torch 2.11.0 and replace CUDA 13.2
packages with CUDA 13.0 packages. That conflicts with the requested latest
cu132 torch stack, so vLLM was installed in \path{/tmp/arena-diffusiongemma-vllm}
for the proof run instead of in the main uv environment. \\
\bottomrule
\end{longtable}

\section{Supported and Unsupported Claims}
\textbf{Supported now.} The course exercise has a working CUDA-trained tiny
discrete diffusion language model, held-out testing, a sampler reconstruction
metric, activation-trajectory checks, and a negative control that fails as
expected. It also has real NVIDIA NVFP4 DiffusionGemma generation evidence in an
isolated vLLM runtime under the 24GB GPU target.

\textbf{Unsupported now.} The course must not claim BF16 direct local inference,
DiffusionGemma denoising-step activation capture, diffusion-time patching,
throughput, broad output quality, or full interpretability parity from this
one-prompt proof. It also must not treat mock generations, fake outputs,
config-only checks, or placeholder artifacts as substitutes for artifact-backed
runs.

\section{Implementation Pointers}
\begin{itemize}
  \item \path{chapter5_modern_architectures/exercises/part5_diffusion_language_models/verification_report.json}
  contains the accepted toy CUDA metrics, BF16 deferral, and NVFP4 isolated-vLLM
  generation proof fields.
  \item \path{chapter5_modern_architectures/exercises/part5_diffusion_language_models/artifacts/diffusiongemma_vllm_probe.json}
  is the compact proof artifact for the real NVFP4 generation run.
  \item \path{arena_ext/diffusion_lm.py} computes the DiffusionGemma readiness
  report and exposes the generation-ready boolean.
  \item \path{arena_ext/gated_artifacts.py} pins the exact Google BF16 and NVIDIA
  NVFP4 repository revisions and required files.
  \item \path{scripts/prepare_diffusiongemma_artifacts.py} inspects or downloads
  the pinned artifacts. The conservative command is:
  \begin{verbatim}
BNB_CUDA_VERSION=130 uv run python scripts/prepare_diffusiongemma_artifacts.py \
  --artifact nvfp4 --download --require-local --max-workers 1
  \end{verbatim}
  \item \path{scripts/run_diffusiongemma_vllm_probe.py} runs the real generation
  probe from an isolated vLLM Python executable.
  \item \path{FOR_REVIEW/tables/diffusiongemma_readiness.csv} is the compact
  evidence table for this packet.
\end{itemize}

\section{Validation Commands}
\begin{verbatim}
uv run python - <<'PY'
import torch, torchvision, modelopt
print(torch.__version__, torch.version.cuda)
print(torch.cuda.get_device_name(0))
props = torch.cuda.get_device_properties(0)
print(round(props.total_memory / 1024**3, 2))
print(torchvision.__version__)
print(modelopt.__version__)
PY

BNB_CUDA_VERSION=130 uv run python -m pytest -q \
  tests/test_diffusion_lm.py \
  chapter5_modern_architectures/exercises/part5_diffusion_language_models/tests.py

BNB_CUDA_VERSION=130 uv run python scripts/run_extension_verification_reports.py \
  --section 1.6 --section 5.1 --section 5.5 --section 5.6 --max-vram-gb 24

BNB_CUDA_VERSION=130 uv run python scripts/prepare_diffusiongemma_artifacts.py \
  --artifact all --json

/tmp/arena-diffusiongemma-vllm/bin/python scripts/run_diffusiongemma_vllm_probe.py \
  --model /home/bitwise/.cache/huggingface/hub/models--nvidia--diffusiongemma-26B-A4B-it-NVFP4/snapshots/2ea837236295d617ac27f8c17d61228081932c40 \
  --prompt 'In one concise sentence, explain why mechanistic interpretability experiments need negative controls.' \
  --max-tokens 128 --max-model-len 4096 --max-num-seqs 1 \
  --gpu-memory-utilization 0.85 --canvas-length 256 --use-chat-template \
  --output-json chapter5_modern_architectures/exercises/part5_diffusion_language_models/artifacts/diffusiongemma_vllm_probe.json

uv pip check
uv pip install --dry-run vllm --prerelease allow
tectonic FOR_REVIEW/FOR_REVIEW.tex --outdir FOR_REVIEW
pdftotext FOR_REVIEW/FOR_REVIEW.pdf -
CHECK=/home/bitwise/.codex/skills/for-review-pdf/scripts/check_for_review_pdf.py
python "$CHECK" FOR_REVIEW/FOR_REVIEW.pdf --tex FOR_REVIEW/FOR_REVIEW.tex
\end{verbatim}

\section{Limitations and Blockers}
\begin{itemize}
  \item The BF16 model is intentionally deferred for local 24GB execution. That
  deferral should be reviewed, but it does not block the intended 24GB NVFP4
  route.
  \item The installed Transformers stack still cannot directly consume the
  NVIDIA ModelOpt NVFP4 quantization config; vLLM is the proven local runtime.
  \item vLLM currently resolves to a stack change that downgrades torch and CUDA
  packages, so it was not installed into the main uv environment. The proof uses
  an isolated environment instead.
  \item The NVFP4 proof is one deterministic prompt. It proves loading and
  generation on this GPU, not output quality, throughput, or a prompt suite.
  \item The exercise still needs released-checkpoint denoising-step
  output/activation capture and a matched random-control diffusion-time patching
  exercise before claiming mechanism-level parity.
\end{itemize}

\section{Questions for Expert Review}
\begin{enumerate}
  \item Is the \status{deferred} marker for direct BF16 local inference accepted
  as a reasonable 24GB hardware deferral?
  \item Is the isolated-vLLM split acceptable for the 24GB local route while the
  main uv environment remains torch 2.12.1+cu132?
  \item Should the next sprint prioritize a ModelOpt-aware Transformers path, a
  documented isolated-vLLM workflow, or a containerized runtime?
  \item What minimum evidence should be required after this one-prompt proof: a
  prompt suite, throughput/VRAM measurements, denoising trajectories, or a
  patching intervention?
  \item What output-quality and reproducibility checks would make the notebook
  feel closest to original ARENA quality rather than a thin model-loading demo?
\end{enumerate}

\appendix
\section{Artifact Index}
\begin{itemize}
  \item \path{Extension-Roadmap.md}: roadmap requirement that quantized
  DiffusionGemma run under the 24GB budget and that every notebook end with a
  verification block.
  \item \path{docs/local_gpu_setup.md}: local setup notes and DiffusionGemma
  artifact preparation command.
  \item \path{docs/method_registry.md}: method registry row marking
  DiffusionGemma local inference as NVFP4-proven in an isolated vLLM runtime.
  \item \path{chapter5_modern_architectures/exercises/part5_diffusion_language_models/artifacts.lock.yml}:
  pinned claim scope, controls, revisions, and expected metrics.
  \item \path{chapter5_modern_architectures/exercises/part5_diffusion_language_models/expected_outputs/reference_metrics.json}:
  expected metric thresholds and current generation-ready flag.
  \item \path{scripts/run_diffusiongemma_vllm_probe.py}: exact probe runner for
  the isolated vLLM generation artifact.
  \item \path{chapter5_modern_architectures/exercises/part5_diffusion_language_models/artifacts/diffusiongemma_vllm_probe.json}:
  prompt, output, vLLM version, torch/CUDA version, GPU name, memory, and timing
  for the real NVFP4 proof.
\end{itemize}

\end{document}
